\bigskip
\begin{center}
{\large\bfseries Stealthy-attack damage bounds for remote estimation under trusted feedback\par}
\end{center}
\noindent\textbf{Abstract:}
This paper studies statistical stealth in stochastic event-triggered remote estimation when transmission events and feedback of the receiver's actual prediction are trusted, but transmitted payloads may be modified after the event decision. For a finite-dimensional linear Gaussian system with noisy full-state measurements and an invertible state matrix, we decompose true mean-square estimation damage into conditional bias energy and a change in the distribution of covariance labels. A normalized next-silence statistic links the hidden bias to the joint law of the trusted event and covariance label. At each fixed finite target time, a uniform upper bound and a feasible causal construction show that the worst-case damage under peak and total-variation constraints has order \(\min\{b^2,\varepsilon\}\). We then consider a scalar protocol in which the actual prediction feedback is quantized and the silent update uses the same quantized center. Under a specified telemetry record, a cell-confined attack can preserve the record law while causing positive damage; the zero-TV worst-case damage has order \(\min\{b^2,\delta^2\}\). Existing calculations and a prespecified paired experiment illustrate the damage decomposition, the analytic two-budget bounds, and the quantized witness. No event split was observed in the finite exact-feedback sample, so the zero count is not treated as evidence of zero population leakage. The quantized result does not establish indistinguishability from the complete receiver log or a fixed-bit implementation. All guarantees are finite-horizon results under the stated information structure.

\smallskip
\noindent\textbf{Key words:} stochastic event-triggering; remote state estimation; statistical stealth; trusted feedback; quantization; estimation damage
